\chapter{Objetivos}
Para asegurar que los objetivos de tu proyecto, trabajo de grado o investigación sean claros, alcanzables y medibles, utiliza el enfoque SMART, el cual define cinco criterios fundamentales que todo buen objetivo debe cumplir: \\[1cm]
S	Específico	¿Qué quiero lograr exactamente? ¿Quién está involucrado? ¿Dónde y cómo se hará?\\[1cm]
M	Medible	¿Cómo sabré que se ha logrado el objetivo? ¿Qué indicadores usaré?\\[1cm]
A	Alcanzable	¿Es posible lograrlo con los recursos y el tiempo disponibles?\\[1cm]
R	Relevante	¿Es importante para mi proyecto, carrera o contexto? ¿Tiene sentido hacerlo?\\[1cm]
T	Temporal	¿En cuánto tiempo debo lograrlo? ¿Cuál es el plazo o la fecha límite?\\[1cm]

 Instrucciones para redactar tus objetivos SMART:\\[1cm]
Empieza con un verbo en infinitivo que denote acción, como: Analizar, Desarrollar, Diseñar, Evaluar, Implementar, Identificar, Probar, Comparar, etc. Revisa la taxonomía de Bloom
\\[1cm]
Enfócate en un solo resultado claro por objetivo. Evita que un mismo objetivo tenga múltiples intenciones o acciones.
\\[1cm]
Verifica que el objetivo sea viable con los recursos, tiempo y contexto del proyecto. No propongas metas inalcanzables.


\section{Objetivo general}
Ejemplos\\[1cm]
Desarrollar un modelo de red neuronal convolucional (CNN) para clasificar imágenes de hojas de plantas en 5 categorías de enfermedades, utilizando un conjunto de datos públicos y validación cruzada.
\\[1cm]
Diseñar un chatbot basado en IA conversacional que resuelva al menos el 70 por ciento de las consultas frecuentes de usuarios en una plataforma educativa, midiendo la tasa de resolución y satisfacción del usuario.
\\[1cm]
Construir un modelo de machine learning para predecir la probabilidad de diabetes tipo 2 en pacientes adultos, utilizando datos clínicos anonimizados

\subsection{Objetivos especificos}
Ejemplos\\[1cm]
Implementar un modelo de aprendizaje automático supervisado basado en BERT para clasificar el sentimiento de 10.000 reseñas de productos en positivo, negativo o neutro.
\\[1cm]
Entrenar un modelo de detección de objetos en tiempo real con YOLOv8 para identificar peatones, señales de tránsito y otros vehículos
\\[1cm]
Entrenar un modelo CNN en TensorFlow para clasificar imágenes de tumores cerebrales en benignos o malignos
\\[1cm]
Comparar el rendimiento de tres arquitecturas de redes neuronales (ResNet, VGG y EfficientNet) usando métricas de exactitud, precisión y recall en un conjunto de imágenes médicas.
\\[1cm]
Desarrollar un modelo basado en BERT para clasificar automáticamente correos de servicio al cliente en cinco categorías, con una precisión mayor al 85 por ciento en el conjunto de prueba.
\\[1cm]
Integrar sensores IoT con una red neuronal para predecir fallas de válvulas industriales, evaluando la tasa de alertas correctas y falsas en un entorno simulado